\documentclass[../main.tex]{subfiles}

\begin{document}
\subsection{{\Struct}}

{\Struct} merupakan kumpulan berbagai tipe data yang berbeda dikelompokkan dalam satu deklarasi unit. Berbeda dengan {\sarray} yang hanya berupa kumpulan variabel bertipe data sama, {\struct} bisa memiliki variabel-variabel yang bertipe data sama atau berbeda, bahkan bisa menyimpan variabel yang bertipe data {\sarray} atau {\struct}  itu sendiri \parencite{Bachtiar_2005}.

\subsection{Deklarasi {\Struct}}

Cara mendeklarasikan sebuah {\struct} adalah harus menggunakan kata kunci \citecode{struct} yang diikuti dengan nama struktur. Kemudian diikuti dengan blok yang diapit dengan tanda kurung kurawal \citecode{\{\}} dimana isinya adalah variabel-variabel (disebut \textit{member} atau \textit{field}) yang akan dideklarasikan di dalam struktur tersebut, serta diakhiri dengan tanta titik koma (\textit{semicolon}) \parencite{utami10}. Bentuk umum deklarasi struct sebagai berikut:

\begin{code}
\begin{verbatim}
struct nama_struktur {
  tipe_data field1;
  tipe_data field2;
  ...
} var1, var2, ...;
\end{verbatim}
\end{code}

\subsection{Akses {\Struct}}

\textit{Struct} dapat diakses dengan menggunakan operator \textit{unary}/operator titik \citecode{.} yang dituliskan di antara nama struktur dan nama \textit{field} \parencite{utami10}. Cara pendeklarasian lainnya yaitu:

\begin{enumerate}
    \item Bila {\struct} dideklarasikan sebagai sebuah variabel biasa, maka cara pengaksesan anggota {\struct} menggunakan operator titik/operator \textit{unary} \citecode{.}.
    \item Bila {\struct} dideklarasikan sebagai sebuah variabel {\pointer}, yaitu cara pengaksesan anggota {\struct} menggunakan operator panah \citecode{→}.
    \item Bila {\struct} dideklarasikan sebagai sebuah variabel dalam {\pointer} sebagai {\sarray}, maka cara pengaksesan anggotanya {\struct} menggunakan operator titik/operator \textit{unary} \citecode{.}.
\end{enumerate}

Contoh \textit{source code struct} berdasarkan pengaksesan sebagai berikut \parencite{munir2016algoritma}:

\subsubsection{{\Struct} dideklarasikan sebagai sebuah variabel biasa}

\begin{indentpar}\hpar{}
    {\Struct} dideklarasikan sebagai sebuah variabel biasa, maka cara pengaksesan anggota {\struct} menggunakan operator titik/operator \textit{unary} \citecode{.}. Contoh \textit{source code}-nya sebagai berikut:
\end{indentpar}

\begin{code}[n]
\begin{verbatim}
#include <iostream>
#include <string.h>
struct personal {
  char nama[20];
  char alamat[20];
  long int telepon;
  char kata[40];
  long int kodepos;
}
void main(){
  //deklarasikan struktur
  personal person:
  //mengisi anggota struktur
  strcpy (person.nama. “Anton”);
  strcpy (person.alamat. “Jl.Majapahit.No.62”);
  serson.telepon=636126;
  strcpy (person.kota, “mataram”);
  person.kodepos=83123;
  cout<< “nama:”<<person.nama<<endl;
  cout<< “alamat:”<<person.alamat<<endl;
  cout<< “telepon:”<<person.telepon<<endl;
  cout<< “kota:”<<person.kota<<endl;
  cout<< “kode pos:”<<person.kodepos<<endl;
}
\end{verbatim}
\end{code}

\subsubsection{{\Struct} dideklarasikan sebagai sebuah variabel {\pointer}}

\begin{indentpar}\hpar{}
    {\Struct} dideklarasikan sebagai sebuah variabel {\pointer}, yaitu cara pengaksesan anggota {\struct} menggunakan operator panah \citecode{→}. Contoh \textit{source code}-nya sebaga berikut:
\end{indentpar}

\begin{code}[n]
\begin{verbatim}
#include <iostream>
#include <stdlib.h>

struct time{
  int jam, min;
};
struct rencana {
  struct time*awal;
  struct timr*akhir;
};
struct time jk={1,2};
struct time ki={3,4};
struct rencana kerja ={&jk, &ki};
void main(){
  kerja.akhir→min > 37;
  cout << kerja.akhir→jam << " " << kerja.awal→min << " " << kerja
  .akhir→jam << " " << kerja.akhir→min << endl;
  system("pause")
}
\end{verbatim}
\end{code}
\vspace{-1ex}
\subsection{{\Nested} {\Struct}}
{\Nested} {\struct} merupakan suatu {\struct} dapat digunakan di dalam {\struct} yang lainnya. Bentuk umum pendeklarasiannya sebagai berikut:

\begin{code}
\begin{verbatim}
struct dtmhs {
	char nim[9];
	char nama[15];
}
struct dtnil{
	float nil1;
	float nil2;
}
struct {
	struct dtmhs mhs;
	struct dtnil nil;
}nilai;begin{lstlisting}
\end{verbatim}
\end{code}

\end{document}
